\section{4교시 (90분) — 실습 ① 변경요청서·영향 분석·CCB 의결서 + EVM·ES 손계산}

% =====================================================================
% 4교시 — 실습 ①
% =====================================================================
\begin{frame}{이 교시에서 하는 것 — 오전 이론 둘을 손으로 쓴다}
\begin{itemize}
\item \textbf{⑪ 변경요청서} — 변경 4건 중 \textbf{CR-09(가맹 발주 마감 20:00 \ra{} 22:00)}를 6축으로 재고 경로를 판정한다
\item \textbf{⑫ 성과보고} — 9/30 누적 PV 77.83 · EV 73.43 · AC 71.13에서 지수·EAC·ES를 \textbf{손으로} 계산한다
\item 두 문서는 서로의 입력 — CR-09의 FP +35가 EAC4의 ETC 행에, A10 $-$10이 CR-09 일정 축에
\item 도구는 종이 캔버스 · 계산기 · 워크북 §1.4·§2.4 빈 템플릿 — 노트북은 검산용
\item 시간 — 안내 10분 / ⑪ 35분 / ⑫ 35분 / 짝 교환 검산 10분
\end{itemize}
\keymsg{정답은 워크북 §1.5·§2.5에 있다 — 보지 말고 쓰고, 다 쓴 뒤에 대조하라}
\note{워크북 §1·§2 실습. Day2 4교시와 같은 진행 — 과제 지시 두 프레임, 작성 요령 네 프레임, 흔한 오류 한 프레임, 예시본 대조, 정리. 오전에 배운 3교시(6축·처리 경로 규칙 ①\til{}⑥)와 1·2교시(편차·지수·EAC 4공식·Earned Schedule)가 이 90분의 재료다. 두 과제를 하나의 교시에 묶은 이유를 먼저 말한다 — 변경요청서의 원가 축(+0.19억)은 성과보고의 상향식 ETC "SI 승인·예정 변경 대가 1.02" 행의 일부이고, 성과보고의 A10 TF 0·$-$10은 변경요청서 일정 축의 전제다. 한쪽만 쓰면 다른 쪽이 틀린다. 시간 배분을 화면에 두고 35분 단위로 종을 울린다. 예상 소요 3분.}
\end{frame}

\begin{frame}{과제 ⑪ — CR-09 가맹 발주 마감 20:00 \ra{} 22:00, 항목 1\til{}8을 직접 쓴다}
\begin{itemize}
\item 접수 — 2026-09-17 협의회 공문 제2026-14호, 요청자 S-10 서정필 협의회장, FZ-10(10-28) 전
\item 기준선 문장 — RTM StRS-073 "접수\til{}출고 확정 12h" \ra{} FRN-FLOW-01 §2 "마감 20:00, 피킹 21:00, 확정 익일 08:00"
\item 대안 넷 — (a) 22:00·야간조 연장 (b) 22:00·14h (c) \textbf{파라미터화, 초기 20:00} (d) 기각
\item 쓸 것 — 6축 영향(항목 4\til{}6) · 허용범위 판정과 처리 경로(항목 7) · CCB 의결서(항목 8)
\item 주어진 것 — 수행사 FP 산정 +35(파라미터 12 + 배치 15 + SAP 시각 8), A10 TF 0 · 현재 $-$10
\end{itemize}
\keymsg{묻는 것은 하나 — 이 변경을 P1 CCB가 결정할 수 있는가, 결정할 수 없는 부분은 어디로 보내는가}
\note{워크북 §1.4 빈 템플릿에 쓴다. CR-08(스폰서 범위 추가)·CR-10(거버넌스)·CR-11(범위 축소)은 예시본으로 읽기만 하고, CR-09만 학생이 쓴다 — 넷 중 유일하게 "P1 권한 밖의 것"이 섞여 있어 처리 경로 판정이 연습이 되기 때문이다. 항목 1\til{}3은 위 정보를 옮겨 적으면 되므로 5분 안에 끝내게 하고, 35분의 대부분을 항목 4\til{}8에 쓰게 한다. 학생이 자주 묻는 것 — "대안을 하나 골라야 하나?" 그렇다. CCB 의결서는 대안 하나에 조건을 붙여 의결하는 문서이고, 대안 넷을 나열만 한 의결서는 의결이 아니다. 추정치가 필요한 곳(R-16 확률 변화, TC 건수)은 "[추정]"을 붙여 적게 한다 — 숫자를 비워 두는 것보다 근거 있는 추정이 낫다. 예상 소요 4분.}
\end{frame}

\begin{frame}{과제 ⑫ — 9/30 누적 세 숫자에서 열두 값을 손으로}
\begin{tbl}[\scriptsize]{L{2.6cm}L{5.4cm}L{3.2cm}}
단계 & 계산할 것 & 입력 \\ \midrule
편차·지수 & SV · CV · SPI · CPI & PV 77.83 · EV 73.43 · AC 71.13 \\
예측 & EAC1 · EAC2 · EAC3 · EAC4 & BAC 145.80 · 상향식 ETC 79.72 \\
목표 효율 & TCPI(BAC) · TCPI(집행 한도 156.00) & BAC · 156.00 \\
Earned Schedule & ES · SV(t) · SPI(t) · IEAC(t) & PV(8월) 71.74 · AT 9 · PD 17 \\
\end{tbl}
\begin{itemize}
\item 소수 셋째 자리까지, 산식을 값 옆에 반드시 쓴다 — 산식 없는 값은 검산 불가
\item 본값을 먼저 끝내고, 시간이 남으면 조정(시차 1.96 · 라이선스 4.70)을 한 줄 더
\end{itemize}
\keymsg{계산기는 써도 된다 — 공식을 외우지 말고, 각 공식의 전제를 값 옆에 한 단어로 적어라}
\note{워크북 §2.4 빈 템플릿 항목 4\til{}6. 월별 PV·EV·AC 표(항목 3)는 이미 채워진 것을 나눠 주고 누적 행만 쓰게 한다 — 이 교시의 목적은 표 작성이 아니라 지수와 예측의 손계산이다. 상향식 ETC 79.72는 항목별 근거표(SI 잔여 38.51, 변경 대가 1.02, 상용SW 5.90, 인프라 7.92, 단말 8.90, 전환 5.10, 시험·감리 7.40, PMO 4.80, 대기 0.18)를 주고 합만 확인하게 한다. 학생이 EAC1이 BAC보다 작게 나오면 당황하는데, 그것이 정답이다 — 141.23이 나오는 이유를 "고정가에서 PMB 아래 EAC는 불가능"이라고 판단하는 것이 2교시의 착시 논의를 손으로 확인하는 지점이다. ES는 PV(8월) 71.74와 PV(9월) 77.83 사이에서 EV 73.43이 어디에 있는지 보간하는 것 — 그림 2-4를 벽에 띄워 둔다. 예상 소요 4분.}
\end{frame}

\begin{frame}{작성 요령 ① — 6축의 각 축에 무엇을 적는가}
\begin{tbl}[\scriptsize]{L{1.4cm}L{5.6cm}L{4.0cm}}
축 & 적는 것 & CR-09에서 \\ \midrule
범위 & WBS 번호 · RTM 변경 이력 · SRS 건수 · FP 델타와 \textbf{누적} & 1.4.2 · 1.5.4 · SRS 4건 · +35 (누적 +245) \\
일정 & 어느 활동의 \textbf{TF}가 얼마 줄었나 — 일수가 아니라 여유 & A10 TF 0 · $-$10에 +1영업일 \\
원가 & FP $\times$ 단가 · 자금원(컨틴전시 풀 / 관리예비비) & 35 $\times$ 0.0055 = +0.19 · 미배정 풀 \\
품질 & TC 재작성·추가 건수 · UAT 수정 · 커버리지 & TC 3건 재작성 · +8 · UAT-FRN-039 \\
리스크 & 어느 R-의 P·EMV가 움직이나 · 잔여 EMV 합 순증 & R-16 0.24 \ra{} 0.36 · R-02 하락 · 순증 +0.05 \\
계약/편익 & 변경 대가 · 편익 ID 유지 여부 · \textbf{타 프로젝트·운영 영향} & +0.19 · B-08 12h 유지 · 물류본부·3PL \\
\end{tbl}
\keymsg{여섯 칸은 모두 "무엇에서 무엇으로"다 — 기준선 문장을 인용하지 않은 축은 빈 칸과 같다}
\note{워크북 §1.3 항목 4\til{}6 작성 가이드. 학생이 가장 잘못 쓰는 축은 일정과 계약/편익이다. 일정 축에 "+1일"이라고 쓰면 아무 뜻이 없다 — 어느 활동인지, 그 활동의 TF가 얼마인지, 임계경로인지가 있어야 허용범위 판정(컷오버 창 0)이 된다. CR-09는 A10이 이미 $-$10이므로 "+1영업일이 허용범위 내"라는 판정은 성과보고의 회복 계획을 전제로 해서만 성립한다 — 그 전제를 판정 칸에 적는 것이 정답이다. 계약/편익 축의 "타 프로젝트·운영 영향" 칸은 CR-09의 핵심이다. 22:00 마감은 P1 화면의 변경이 아니라 이천센터 야간 작업과 대성로지스 배송 출발 시각의 변경이며, 그 칸이 있어야 처리 경로 규칙 ④(운영위 보고 병행)가 발동한다. 범위 축의 누적 +245는 CR-08까지의 +210에 +35를 더한 것 — 허용범위 ±372의 66\%다. 예상 소요 5분.}
\end{frame}

\begin{frame}{작성 요령 ② — 처리 경로 판정과 CCB 의결서의 네 칸}
\begin{itemize}
\item 판정 순서 — 6축 허용범위 \ra{} 헌장 문언 변경 여부 \ra{} 타 프로젝트·운영 영향 \ra{} 관리예비비 필요
\item CR-09 — 6축 모두 내 · 헌장 문언 불변 · \textbf{운영 영향 예} \ra{} 규칙 ④ · PM 전결 불가
\item 경로 — 변경협의체(09-30 합의) \ra{} CCB(10-15) + 프로그램 운영위 보고(10-08)
\item 의결서 네 칸 — \textbf{의결}(대안 하나) · \textbf{조건}(누가 언제까지) · \textbf{반대의견}(서면 그대로) · \textbf{미이행 시}
\item 설계 변경과 운영 변경을 \textbf{분리}한다 — 파라미터(P1 권한)는 승인, 22:00 운용(물류 권한)은 재상정
\end{itemize}
\keymsg{CCB가 자기 권한 밖의 것을 결정하면 집행되지 않고, 미루면 요구가 방치된다 — 분리해서 의결하라}
\note{워크북 §1.3 항목 7·8과 §1.5 처리 경로 규칙 ①\til{}⑥. 판정 순서를 화면 순서대로 고정시킨다 — 학생들은 흔히 "FP +35는 작으니 PM 전결"로 끝내는데, 허용범위 내라는 것은 규칙 ①의 첫 조건일 뿐이고 규칙 ④(운영 영향)가 걸리면 CCB다. 의결서의 정답 골격 — 조건부 승인 · 대안 (c) · FP +35 · 0.19억 · 초기 운용 20:00. 조건 ① 22:00 전환은 물류본부 야간조 편성안과 3PL 배송 출발 시각 합의를 첨부해 2027-01 CCB에 재상정, ② 협의회 회신 10-23, ③ B-08 12h 유지·14h 불채택, ④ R-16 잔여 EMV 갱신. 미이행 시 — 2027-01 CCB 미상정이면 20:00 운용 확정. 반대의견 칸에 한동석의 서면("야간조 편성은 P2 안정화 이후, 파라미터 구현에는 이의 없음")을 그대로 적는 것이 의결서를 나중에 읽는 사람을 위한 최소한이다. 이해관계자 퇴장(이재현, 변경 대가 확인 후 의결 시)도 기록 항목이다. 예상 소요 5분.}
\end{frame}

\begin{frame}{작성 요령 ③ — EVM 손계산의 순서: 세 숫자에서 열두 값으로}
\begin{itemize}
\item 입력 셋 — PV 77.83 · EV 73.43 · AC 71.13 (누적, 2026-09-30) · BAC 145.80
\item 편차 둘 — SV = EV $-$ PV = $-$4.40 · CV = EV $-$ AC = +2.30 (부호를 먼저 적는다)
\item 지수 둘 — SPI = EV/PV = 0.943 · CPI = EV/AC = 1.032 (소수 셋째 자리)
\item EAC 넷 — BAC/CPI = 141.23 · AC+(BAC$-$EV) = 143.50 · AC+(BAC$-$EV)/(CPI×SPI) = 145.43 · AC+상향식 ETC 79.72 = 150.86
\item TCPI 둘 + VAC — TCPI(BAC) = (BAC$-$EV)/(BAC$-$AC) = 0.969 · TCPI(156.0) = 0.853 · VAC = BAC $-$ EAC (넷 다)
\end{itemize}
\keymsg{순서를 바꾸지 마라 — EAC를 먼저 쓰고 CPI를 맞추는 학생이 매 기수 있다. 공식은 전제의 이름이지 결과의 이름이 아니다}
\note{워크북 §2.3 항목 3\til{}5, §2.6 검산식(src/day3\_evm.py). 판서 순서를 이 다섯 줄 그대로 고정한다. 학생이 가장 자주 틀리는 자리 셋 — (1) SV 부호(PV$-$EV로 계산해 +4.40), (2) EAC3의 분모(CPI×SPI를 CPI+SPI로), (3) TCPI 분모에 BAC 대신 EAC를 쓰고도 "TCPI(BAC)"라고 적는 것. 상향식 ETC 79.72의 내역(SI 잔여 38.51 + 변경 1.02 + 대기 0.18 + 상용SW 5.90 + 인프라 7.92 + 단말 8.90 + 전환 5.10 + 시험·감리 7.40 + PMO 4.80)은 워크북 표를 옆에 두고 옮겨 적게 한다 — 이것이 "공식이 아니라 근거"라는 2교시 명제의 실습판. 착시 조정(AC 75.83 \ra{} CPI 0.968)은 심화 과제. 예상 소요 15분(학생 작업 포함).}
\end{frame}

\begin{frame}{작성 요령 ④ — Earned Schedule 읽기: EV가 PV 곡선에서 "언제"인가}
\begin{itemize}
\item ES = 8 + (EV $-$ PV$_8$) / (PV$_9$ $-$ PV$_8$) = 8 + (73.43 $-$ 71.74) / (77.83 $-$ 71.74) = \textbf{8.28개월}
\item AT = 9 (데이터 일자 2026-09-30 = 착수 후 9개월)
\item SV(t) = ES $-$ AT = $-$0.72개월 (≈ $-$15영업일) · SPI(t) = ES/AT = \textbf{0.920}
\item IEAC(t) = PD / SPI(t) = 17 / 0.920 = \textbf{18.48개월} \ra{} 추세 종료 2027-07-12 (허용 G4+15 = 06-18)
\item 대조 — SPI 0.943(돈) vs SPI(t) 0.920(시간): 7월부터 벌어지고 9월 차이 $-$0.023
\end{itemize}
\keymsg{ES는 "EV 73.43을 PV 곡선이 언제 지났는가"를 자로 재는 것 — 그림 2-4의 빨간 점을 손으로 찾으면 끝난다}
\note{워크북 §2.3 항목 6, 그림 2-4. 계산보다 읽기가 먼저 — 그림 2-4를 배포하고 EV 73.43을 PV 곡선에서 찾아 x축을 읽게 한다(8월과 9월 사이, 8월에 더 가까움). 그다음 보간식으로 8.28을 확인한다. PV$_8$ = 71.74·PV$_9$ = 77.83은 Day2 §4 월별 PV 표의 정본. IEAC(t)를 날짜로 옮길 때 17개월 PD의 기준일(2026-01-05)에서 18.48개월을 더하면 2027-07-12 — 허용범위 종료 06-18보다 늦으므로 일정 축 예외(EX-01)의 정량 근거가 된다. SPI와 SPI(t)의 차이가 왜 7월부터 벌어지는지(PV 곡선이 5\til{}6월 조달 봉우리 뒤로 완만해져 같은 EV 부족이 더 긴 시간으로 환산됨)는 심화 질문. 예상 소요 12분(학생 작업 포함).}
\end{frame}

\begin{frame}{흔한 오류 — 이 실습에서 매 기수 나오는 다섯}
\begin{itemize}
\item 구두 합의를 변경으로 쓰지 않는다 — "서정필 회장과 이야기됐다"는 CR-09의 접수 근거가 아니라 접수 사유다
\item SPI를 일수로 읽는다 — 0.943 × 9개월 ≠ 지연 일수. PV 곡선에 봉우리가 있으면 틀린다(ES를 쓰라)
\item 월별 지수와 누적 지수를 섞는다 — 4월 누적 SPI 0.784와 4월 월별 EV/PV는 다른 숫자
\item 라이선스 이연을 원가 절감으로 읽는다 — CV +2.30의 +4.70은 아직 안 낸 돈이다
\item 6축 중 "편익" 축을 비운다 — CR-09는 편익 B-08(리드타임 12h)에 직접 걸리는 변경이다
\end{itemize}
\keymsg{다섯 오류의 공통점 — 숫자를 계산은 했는데 그 숫자가 무엇의 이름인지 적지 않았다}
\note{워크북 §1.7·§2.7 점검 질문과 Day1\til{}Day2 실습에서 반복 관찰된 오류를 모았다. 각 조 산출물을 걷기 전에 이 다섯을 화면에 띄우고 자기 것을 스스로 고치게 한다(10분). 특히 4번 — 5월 CPI 1.073 점프를 "원가 효율 개선"으로 적은 조가 있으면 그 조의 성과보고서를 예시로(익명) 띄워 착시 1을 다시 설명한다. 예상 소요 8분.}
\end{frame}

\begin{frame}{예시본 참조 · 시간 배분 · 심화 과제}
\begin{itemize}
\item 예시본 — 워크북 §1.5(CR-09 완성본 + CCB 10-15 의결서) · §2.5(성과보고 제9호, Highlight Report 1장)
\item 인쇄용 — \texttt{../템플릿/PM\_Day3/11\_변경요청서\_완성예시.pdf} · \texttt{12\_성과보고\_완성예시.pdf} (빈 양식은 \texttt{\_빈템플릿.pdf})
\item 시간 — ⑪ 변경요청서·6축·의결서 40분 / ⑫ EVM·ES 손계산 35분 / 자기 점검·대조 15분
\item 자기 점검 — 워크북 §1.7·§2.7 점검 질문 각 3개. 예시본과 값이 다르면 \textbf{어느 전제가 달랐는지}를 적는다
\item 심화(3\til{}7년차) — 자기 프로젝트의 최근 변경 1건을 6축으로 다시 재고, 마지막 성과보고에 SPI(t)를 추가해 본다
\end{itemize}
\keymsg{예시본은 정답이 아니라 전제가 적힌 답이다 — 다른 값이 나왔으면 전제를 비교하라}
\note{Day2 4교시와 같은 운영. 예시본 PDF는 실습 시작 시 배포하지 않고 35분 뒤(⑫ 착수 시점)에 ⑪ 예시본만 먼저 배포, ⑫ 예시본은 손계산이 끝난 뒤. 값이 다른 조가 나오면 "전제"를 찾게 하는 것이 이 실습의 목적이다 — EAC3에서 CPI×SPI에 조정 CPI 0.968을 쓴 조는 155.04(경고 대역)가 나오고, 그것은 틀린 것이 아니라 다른 전제다. 예상 소요 90분 전체.}
\end{frame}

\begin{frame}{4교시 핵심 정리}
\begin{enumerate}
\item CR-09는 FP +35라 작아 보이지만 운영 영향(규칙 ④)으로 CCB — 크기와 권한은 다른 축이다
\item 의결서 네 칸(의결·조건·반대의견·미이행 시) 중 반대의견 칸이 나중에 읽는 사람을 위한 최소한
\item 세 숫자(PV·EV·AC)에서 열두 값 — 순서가 곧 전제. EAC 넷은 네 개의 미래다
\item ES 8.28은 계산이 아니라 읽기 — 그림에서 찾고 식으로 확인한다
\item 다섯 오류의 공통점 — 계산은 했는데 이름을 안 적었다
\end{enumerate}
\keymsg{5교시는 나머지 셋(⑬ 이슈·리스크 / ⑭ 품질·감리 / ⑮ 조달·계약) — 숫자가 아니라 기록이 무기가 되는 영역}
\note{조별 산출물 회수. 성과보고서 1장(Highlight Report)은 6교시 시작 전까지 벽에 붙여 두고 갤러리워크로 서로 대조하게 한다 — RAG 색이 다른 조가 있으면 그것이 5교시 "RAG의 정치" 도입 재료다. 예상 소요 3분.}
\end{frame}
